\documentclass[
	12pt,
%	%BCOR=5mm,
	DIV=14,
	headinclude=on,
	footinclude=off,
	parskip=half,
	bibliography=totoc,
	listof=entryprefix,
	toc=listof,
	pointlessnumbers,
	plainfootsepline]{scrreprt}
	
	%\usepackage[paperwidth=210mm, paperheight=245mm, top=25mm, bottom=20mm, left=25mm,	right=25mm,]{geometry}

%\usepackage[onehalfspacing]{setspace}

% Theorem

%	Konfigurationsdatei einziehen
% !TEX root =  master.tex
% 		HYPERREF
%
\usepackage[
hidelinks=true % keine roten Markierungen bei Links
]{hyperref}

% \newcommand*{\Kapitel}{../Kapitel}


\newcommand{\TitelDerArbeit}[1]{\def\DerTitelDerArbeit{#1}\hypersetup{pdftitle={#1}}}
\newcommand{\AutorDerArbeit}[1]{\def\DerAutorDerArbeit{#1}\hypersetup{pdfauthor={#1}}}
\newcommand{\ZweiterAutorDerArbeit}[1]{\def\DerZweiteAutorDerArbeit{#1}\hypersetup{pdfauthor={#1}}}
\newcommand{\DritterAutorDerArbeit}[1]{\def\DerDritteAutorDerArbeit{#1}\hypersetup{pdfauthor={#1}}}
\newcommand{\Firma}[1]{\def\DerNameDerFirma{#1}}
\newcommand{\Kurs}[1]{\def\DieKursbezeichnung{#1}}



%		FONT AND INPUT ENCODING
%
\usepackage[T1]{fontenc}
\usepackage[utf8]{inputenc}

%		CALCULATIONS
%
\usepackage{calc} % Used for extra space below footsepline

%		LANGUAGE SETTINGS
%
%\usepackage[ngerman]{babel} 	% German language
%\usepackage[german=quotes]{csquotes} 	% correct quotes using \enquote{}

\usepackage[english]{babel}   % For english language
\usepackage{csquotes} 	% Richtiges Setzen der Anführungszeichen mit \enquote{}


%		BIBLIOGRAPHY SETTINGS
%
%\usepackage[backend=biber, autocite=footnote, style=authortitle-ibid, dashed=false,ibidtracker=context]{biblatex} 	%Use Author-Year-Cites with footnotes
 \usepackage[backend=biber, autocite=inline, style=ieee, , sorting=nty, dashed=false]{biblatex} 	% Use IEEE-Style (e.g. [1])
% \usepackage[backend=biber, autocite=inline, style=alphabetic]{biblatex} 	% Use alphabetic style (e.g. [TGK12])
%%%% APA/Harvard-Style (bitte die nächten zwei Zeilen auskommentieren)
%\usepackage[backend=biber, style=apa]{biblatex} 	
%\DeclareLanguageMapping{german}{german-apa}

\DefineBibliographyStrings{ngerman}{  %Change u.a. to et al. (german only!)
	andothers = {{et\,al\adddot}},
}
\DefineBibliographyStrings{ngerman}{%
	urlseen = {Last accessed},
}

%%% Uncomment the following lines to support hard URL breaks in bibliography 
%\apptocmd{\UrlBreaks}{\do\f\do\m}{}{}
%\setcounter{biburllcpenalty}{9000}% Kleinbuchstaben
%\setcounter{biburlucpenalty}{9000}% Großbuchstaben


\setlength{\bibparsep}{\parskip}		%add some space between biblatex entries in the bibliography
\addbibresource{bibliography.bib}	%Add file bibliography.bib as biblatex resource


%		FOOTNOTES 
%
% Count footnotes over chapters
\usepackage{chngcntr}
\counterwithout{footnote}{chapter}


%	ACRONYMS
%%%
%%% WICHTIG: Installieren Sie das neueste Acronyms-Paket!!!
%%%
\makeatletter
\usepackage[nolist]{acronym}
\@ifpackagelater{acronym}{2015/03/20}
{%
	\renewcommand*{\aclabelfont}[1]{\textbf{\textsf{\acsfont{#1}}}}
}%
{%
}%
\makeatother

%		LISTINGS
\usepackage{listings}	%Format Listings properly
\usepackage{xcolor}
\usepackage{bera}% optional: just to have a nice mono-spaced font
%\renewcommand{\lstlistingname}{Quelltext} 
\renewcommand{\lstlistlistingname}{Quelltextverzeichnis}

\lstset{numbers=left,
	numberstyle=\tiny,
	captionpos=b,
	basicstyle=\ttfamily\small}

\colorlet{punct}{red!60!black}
\definecolor{background}{HTML}{EEEEEE}
\definecolor{delim}{RGB}{20,105,176}
\colorlet{numb}{magenta!60!black}

\lstdefinelanguage{json}{
	basicstyle=\normalfont\ttfamily,
	numbers=left,
	numberstyle=\scriptsize,
	stepnumber=1,
	numbersep=8pt,
	showstringspaces=false,
	breaklines=true,
	frame=lines,
	backgroundcolor=\color{background},
	literate=
	*{0}{{{\color{numb}0}}}{1}
	{1}{{{\color{numb}1}}}{1}
	{2}{{{\color{numb}2}}}{1}
	{3}{{{\color{numb}3}}}{1}
	{4}{{{\color{numb}4}}}{1}
	{5}{{{\color{numb}5}}}{1}
	{6}{{{\color{numb}6}}}{1}
	{7}{{{\color{numb}7}}}{1}
	{8}{{{\color{numb}8}}}{1}
	{9}{{{\color{numb}9}}}{1}
	{:}{{{\color{punct}{:}}}}{1}
	{,}{{{\color{punct}{,}}}}{1}
	{\{}{{{\color{delim}{\{}}}}{1}
	{\}}{{{\color{delim}{\}}}}}{1}
	{[}{{{\color{delim}{[}}}}{1}
	{]}{{{\color{delim}{]}}}}{1},
}

%		EXTRA PACKAGES
\usepackage{lipsum}    %Blindtext
\usepackage{graphicx} % use various graphics formats
\usepackage[german]{varioref} 	% nicer references \vref
\usepackage{caption}	%better Captions
\captionsetup[table]{position=bottom} 
\usepackage{booktabs} %nicer Tabs
\usepackage{array}
\usepackage{pdfpages} 
\usepackage{algpseudocode}% http://ctan.org/pkg/algorithmicx
%\newcolumntype{P}[1]{>{\raggedright\arraybackslash}p{#1}}
\usepackage{tikz}
\usepackage{amssymb}

\usepackage{multirow}

\newcommand\MyBox[2]{
	\fbox{\lower0.75cm
		\vbox to 1.7cm{\vfil
			\hbox to 1.7cm{\hfil\parbox{1.4cm}{#1\\#2}\hfil}
			\vfil}%
	}%
}



%chapter settings
\usepackage{etoolbox}

%		ALGORITHMS
\usepackage{algorithm}
\usepackage{amsmath}
\usepackage{algpseudocode}
\usepackage{soul}
\usepackage{bchart}
\renewcommand{\listalgorithmname}{Algorithmenverzeichnis}
\floatname{algorithm}{Algorithmus}


%		FONT SELECTION: Entweder Latin Modern oder Times / Helvetica
\usepackage{lmodern} %Latin modern font
\usepackage{tabularx}
%\usepackage{mathptmx}  %Helvetica / Times New Roman fonts (2 lines)
%\usepackage[scaled=.92]{helvet} %Helvetica / Times New Roman fonts (2 lines)

%		PAGE HEADER / FOOTER
%	    Warning: There are some redefinitions throughout the master.tex-file!  DON'T CHANGE THESE REDEFINITIONS!
\RequirePackage{scrlfile}
\ReplacePackage{scrpage2}{scrlayer-scrpage}
\RequirePackage[automark,headsepline,footsepline]{scrpage2}
\pagestyle{scrheadings}
\renewcommand*{\pnumfont}{\upshape\sffamily}
\renewcommand*{\headfont}{\upshape\sffamily}
\renewcommand*{\footfont}{\upshape\sffamily}
\renewcommand{\chaptermarkformat}{}

\clearscrheadfoot

\ifoot[\rule{0pt}{\ht\strutbox+\dp\strutbox}Advanced Machine Learning]{\rule{0pt}{\ht\strutbox+\dp\strutbox}Advanced Machine Learning}
\ofoot[\rule{0pt}{\ht\strutbox+\dp\strutbox}\pagemark]{\rule{0pt}{\ht\strutbox+\dp\strutbox}\pagemark}

\ohead{\headmark}

\usepackage{chngcntr}
\counterwithout{figure}{chapter}
\counterwithout{table}{chapter}

%zeilenumbrüche
\setlength{\emergencystretch}{1em}

\usetikzlibrary{calc,trees,positioning,arrows,chains,shapes.geometric,%
	decorations.pathreplacing,decorations.pathmorphing,shapes,%
	matrix,shapes.symbols,quotes,angles,babel,calc}
\usepackage{tikz-qtree-compat,tikz-qtree}

\nocite{*}

\tikzset{
	>=stealth',
	punktchain/.style={
		rectangle, 
		rounded corners, 
		% fill=black!10,
		draw=black, very thick,
		text width=10em, 
		minimum height=3em, 
		text centered, 
		on chain},
	line/.style={draw, thick, <-},
	element/.style={
		tape,
		top color=white,
		bottom color=blue!50!black!60!,
		minimum width=8em,
		draw=blue!40!black!90, very thick,
		text width=10em, 
		minimum height=3.5em, 
		text centered, 
		on chain},
	every join/.style={->, thick,shorten >=1pt},
	decoration={brace},
	tuborg/.style={decorate},
	tubnode/.style={midway, right=2pt},
}


\definecolor{myblue}{RGB}{209, 216, 224}

\usepackage{amsthm}

\theoremstyle{plain}
\newtheorem{thm}{Theorem}[chapter] % reset theorem numbering for each chapter

\theoremstyle{definition}
\newtheorem{defn}[thm]{Definition} % definition numbers are dependent on theorem numbers
\newtheorem{exmp}[thm]{Example} % same for example numbers

\usepackage{epstopdf}

\makeatletter
\newcommand{\chapterauthor}[1]{%
	{\parindent0pt\vspace*{-35pt}%
		\linespread{1.1}\large\scshape#1%
		\par\nobreak\vspace*{5pt}}
	\@afterheading%
}
\makeatother

\makeatletter
\newcommand{\sectionauthor}[1]{%
	{\parindent0pt\vspace*{-20pt}%
		\linespread{1.1}\large\scshape#1%
		\par\nobreak\vspace*{5pt}}
	\@afterheading%
}
\makeatother

\makeatletter
\newcommand{\subsubsectionauthor}[1]{%
	{\parindent0pt\vspace*{-20pt}%
		\linespread{1.1}\scshape#1%
		\par\nobreak\vspace*{5pt}}
	\@afterheading%
}
\makeatother

\addbibresource{bibliography.bib}


\begin{document}

% Disable listing with chapter numbers
\counterwithout{lstlisting}{chapter}
\counterwithout{equation}{chapter}
% CHAPTER CHANGES
\renewcommand*{\chapterheadstartvskip}{\vspace*{0cm}}
\renewcommand*{\chapterheadendvskip}{\vspace{2cm}}


\makeatletter
\patchcmd{\scr@startchapter}{\if@openright\cleardoublepage\else\clearpage\fi}{}{}{}
\makeatother

\TitelDerArbeit{SIIM-FISABIO-RSNA COVID-19 Detection}
\AutorDerArbeit{Tobias Richstein (3596554)}
\ZweiterAutorDerArbeit{Julian Seibel (3601340)}
\DritterAutorDerArbeit{Torben Krieger (MatNr)}
\begin{titlepage}
\begin{minipage}{\textwidth}
		\vspace{-2cm}
		\centering
		 \includegraphics[scale=0.35]{img/uni_blacksmall.png}
\end{minipage}
\vspace{6em}
\sffamily
\begin{center}
	\textsf{\large{}Ruprecht-Karls-Universität \\ Heidelberg}\\[2em]
	\textsf{\textbf{\Large{}Final project report}}\\[3mm]
	\textsf{\textbf{Advanced Machine Learning}} \\[3mm] \textsf{Prof. Dr. Köthe} \\[1.5cm]
	\textsf{\textbf{\Large{}\DerTitelDerArbeit}}\\[1.5cm]

	
	\vspace{3em}
	%\textsf{\Large{Sperrvermerk}}
\vfill

\begin{minipage}{\textwidth}

\begin{tabbing}
	Wissenschaftlicher Betreuer: \hspace{0.85cm}\=\kill
	Authors: \> \DerAutorDerArbeit \\ \> \DerZweiteAutorDerArbeit \\ \>  \DerDritteAutorDerArbeit \\ [1.5mm]
	Field of Studies: \> M. Sc. Applied Computer Science \\[1.5mm]
	Period: \> Summer term 2021
\end{tabbing}
\end{minipage}

\end{center}

\end{titlepage}

\pagenumbering{Roman} % Römische Seitennummerierung
\normalfont

%--------------------------------
% Verzeichnisse - nicht benötige Verzeichnisse bitte auskommentieren / löschen.
%--------------------------------

%Sperrvermerk
%\input{Kapitel/sperrvermerk.tex}

%\input{Kapitel/widmung.tex}

%	Kurzfassung
%\chapter*{Kurzfassung}
\begingroup
\begin{table}[h!]
\setlength\tabcolsep{0pt}
\begin{tabular}{p{3.7cm}p{11.7cm}}
Titel: & \DerTitelDerArbeit \\
Verfasser: & \DerAutorDerArbeit \\
\end{tabular}
\end{table}
\endgroup



%	Inhaltsverzeichnis
\tableofcontents

%	Abbildungsverzeichnis
%\listoffigures

%	Tabellenverzeichnis
%\listoftables

%	Listingsverzeichnis
 %\lstlistoflistings

% 	Algorithmenverzeichnis
 %\listofalgorithms

% 	Abkürzungsverzeichnis (siehe Datei acronyms.tex!)

\begin{acronym}[]
	\acro{CNN}{Convolutional Neural Network}
	\acro{IoU}{Intersection over Union}
	\acro{mAP}{mean Average Precision}
	\acro{R-CNN}{Region Based CNN}
	\acro{SGD}{Stochastic Gradient Descent}
	\acro{CAD}{Computer Aided detection and diagnosis}
	\acro{CIoU}{Complete-IoU}
	\acro{SVM}{Support Vector Machine}
	\acro{YOLO}{You only look once}
	\acro{RT-PCR}{Real‐Time reverse transcription Polymerase Chain Reaction}
	\acro{CT}{Computed Tomography}
	\acro{MSE}{mean squared error}
	\acro{GIoU}{Generalized Intersection over Union}
	\acro{DIoU}{Distance Intersection over Union}
	\acro{NIH}{National Institutes of Health}
	\acro{RSNA}{Radiological Society of North America}
	\acro{CXR}{Chest X-Ray}
	\acro{PNG}{Portable Network Graphics}
	\acro{COCO}{Common Objects in Context}
\end{acronym}


%--------------------------------
% Start des Textteils der Arbeit
%--------------------------------
\clearpage
\ihead{\chaptername~\thechapter} % Neue Header-Definition
\pagenumbering{arabic}  % Arabische Seitenzahlen

%	Anleitungs-Datei anleitung.tex einziehen. Auf diese Weise sollten Sie versuchen, für jedes einzelne
% Kapitel eine eigene Datei anzulegen und mittels input-Kommando einzuziehen.
%% !TEX root =  master.tex
\chapter{Gebrauchsanleitung}

\section{Übersicht über die Vorlage}
Die Vorlage wurde im UTF-8 Encoding erstellt. Sollten daher z.\,B. Umlaute in Ihrem \LaTeX-Editor nicht korrekt dargestellt werden, überprüfen Sie bitte die Encoding-Einstellungen des Editors. In seltenen Fällen müssen Sie die Vorlage danach noch einmal neu in den Editor einbinden. 
Die Vorlage beinhaltet die folgenden, in Tabelle \vref{tab:dateien} aufgelisteten Dateien: 
\begin{table}[h!]
	\centering
\begin{tabular}{lp{10cm}}
	\textbf{Dateiname} & \textbf{Beschreibung}\\\toprule
	\texttt{master.tex} & Die Hauptdatei. Alle anderen Dateien werden von dieser Datei eingezogen. \\
	\texttt{abstract.tex} & Die Kurzfassung der Arbeit. \\	
	\texttt{config.tex} & Konfigurationseinstellungen 	 der einzelnen Pakete\\
	\texttt{acronyms.tex} & Definition von Abkürzungen. \\
	\texttt{titlepage.tex} & Titelseite der Arbeit. \textbf{Bitte Anpassen!}\\
	\texttt{anleitung.tex} & Diese Anleitung\\ 
	\texttt{bibliography.bib}&  Die Literaturdatenbank -- hier können Sie die verwendete Literatur einpflegen.\\
	\texttt{ewerkl.tex} & Ehrenwörtliche Erklärung. \textbf{Bitte Anpassen!}\\
	\texttt{appendix.tex} & Anhang bzw. Anhänge \\\bottomrule
\end{tabular}
\caption{\label{tab:dateien}Übersicht über die Dateien der Vorlage}
\end{table}

Es werden -- unter anderem -- die folgenden Zusatzpakete von dieser Vorlage eingezogen und sollten daher in aktuellen Versionen installiert sein: 
\begin{itemize}
	\item\texttt{KOMA-Script} bzw. die Dokumentenklasse \texttt{scrreprt}
	\item\texttt{hyperref} für PDF-Informationen und Links 
	\item \texttt{babel} für länderspezifische Einstellungen
	\item \texttt{csquotes} für sprachabhängige Anführungszeichen (Befehl: \texttt{\textbackslash enquote})
	\item \texttt{acronym} für das Erstellen des Abkürzungsverzeichnisses 
	\item \texttt{booktabs} für das typografisch schöne Setzen von Tabellen 
	\item \texttt{varioref} für einfaches Referenzieren 
	\item \texttt{listings} für schöne Quelltexte
	\item \texttt{algorithm} für schöne Algorithmen
	\item \texttt{bibltatex} und \texttt{biber} für die Erstellung des Literaturverzeichnisses.
\end{itemize}
Alle Konfigurationen dieser Vorlage können in der Datei \texttt{config.tex} eingesehen und ggf. verändert werden. Bitte schauen Sie sich die entsprechenden Dokumentationen 
der Pakete an (\url{https://www.ctan.org}), um deren Verwendung und Möglichkeiten jenseits der hier gezeigten Beispiele zu erlernen.


\section{Übersetzung von \LaTeX-Dateien}
Die Übersetzung von \LaTeX-Dateien erfolgt in mehreren Schritten und unter der Zuhilfenahme unterschiedlicher Programme. Das Hauptdokument (hier die Datei \texttt{master.tex}) wird mittels \texttt{pdflatex} zu einem PDF übersetzt. Ggf. ist eine mehrfache Übersetzung notwendig, um z.\,B. das Inhaltsverzeichnis korrekt darzustellen. 

Für die Einbindung des Literaturverzeichnisses wird nicht mehr das ältere \texttt{bibtex}, sondern das neuere \texttt{biber} in Kombination mit \texttt{biblatex} verwendet. Bitte stellen Sie Ihren \LaTeX-Editor so ein, dass die Verwendung von Biber beim Übersetzungsprozess erfolgt. 

\section{Verwendung von Akronymen}
Akronyme müssen in der Datei \texttt{acronyms.tex} definiert werden (schauen Sie sich hierzu bitte die entsprechende Paket-Dokumentation an!)
Ein definiertes Akronym kann dann mit dem Befehl \texttt{\textbackslash ac} verwenden, so wird z.\,B. \texttt{\textbackslash ac\{DHBW\}} zu \ac{DHBW}. Im weiteren Verlauf wird das 
Acronym dann nur noch in der Kurzform dargestellt: \ac{DHBW}. Die Aufnahme eines verwendeten Akronyms in das Abkürzungsverzeichnis erfolgt automatisch.

\section{Zitieren von Quellen}
Mit dem Befehl \texttt{\textbackslash autocite} kann zitiert werden, z.\,B. so. Sollen mehrere Referenzen auf einmal gesetzt werden, können Sie dies mit dem Befehl \texttt{\textbackslash autocites} erreichen, z.\,B. so. Wird \texttt{autocite} konsequent 
verwendet, kann in der Datei \texttt{config.tex} der Zitationsstil umgeschaltet werden, ohne dass im Text Veränderungen vorgenommen werden müssen. Vorkonfigurierte Stile sind Alphabetic, Harvard, Fußnotenzitat und IEEE-Style. Die Übernahme der Quellen in das Literaturverzeichnis erfolgt automatisch. Ein Beispiel für eine Online-Quelle ist ebenfalls enthalten.\autocite[Vgl.][]{TestOnlineQuelle}

\medskip
Jan stinkt \autocite{Stutenbaumer.2015}


Soll einer Abbildung eine Quellenangabe zugefügt werden, bietet es sich an, diese direkt in der jeweiligen Abbildungsbeschriftung zu hinterlegen. Hierfür kann der Befehl \texttt{\textbackslash cite} verwendet werden, um eine ungewollte Fußnote zu vermeiden. Ein Beispiel ist in Abbildung 
\vref{fig:test} zu sehen. 


\section{Text in Anführungszeichen}
Soll ein Text in Anführungszeichen gesetzt werden, kann dies über den Befehl \texttt{\textbackslash enquote} \enquote{so erreicht werden}. Die Anführungszeichen ändern sich automatisch auf die 
jeweiligen Länderspezifika, wenn die Spracheinstellung des \texttt{babel}-Pakets geändert wird. Voreinstellung ist die deutsche Verwendung von 
Anführungszeichen.




\section{Beispiele}
\lipsum[1]

\subsection{Unterabschnitte}
Es gibt neben \texttt{\textbackslash chapter} auch noch  \texttt{\textbackslash section}, \texttt{\textbackslash subsection}, \texttt{\textbackslash subsubsection} etc. Eine zu starke Untergliederung des Textes sollte jedoch vermieden werden (z.\,B. ein Abschnitt 3.4.2.5.3). 

\subsection{Tabellen und Abbildungen}
Tabellen und Abbildungen sind sogenannte \textit{Floating Objects}, d.\,h. \LaTeX\ setzt diese Objekte an Positionen, die satztechnisch geeignet sind. Daher kann es vorkommen, dass Tabellen oder Abbildungen auf einer anderen Seite erscheinen, die dann referenziert werden müssen. Hier ein Beispiel dafür: 

In Tabelle \vref{tab:tabelle1} ist eine Tabelle abgebildet, die mit dem Befehl \texttt{\textbackslash vref} referenziert wurde. Gleiches kann man auch mit Abbildungen 
machen, wie z.\,B. mit der Abbildung \vref{fig:test}. \LaTeX~ kümmert sich darum, wo die Abbildungen gesetzt werden und passt den Text der Referenz entsprechend an. Soll nur die Nummerierung in den Text geschrieben werden, dann kann auch der Befehl \texttt{\textbackslash ref} verwendet werden.
Abbildungen sollten -- falls möglich -- als Vektor-PDF eingebunden 
werden, da die diese dann beliebig skalieren können.

\lipsum[1]
\begin{table}
	\centering
	\begin{tabular}{p{3cm}crl}
		\textbf{Spalte 1} & \textbf{Spalte 2} & \textbf{Spalte 3} & \textbf{Spalte 4}\\\toprule
		Zeile 1 Spalte 1 &  Zeile 1 Spalte 2 & Zeile 1 Spalte 3 & Zeile 1 Spalte 4\\
		Zeile 2 Spalte 2 &  Zeile 2 Spalte 2 & Zeile 2 Spalte 3 & Zeile 2 Spalte 4\\\midrule
		Zeile 3 Spalte 1 &  Zeile 3 Spalte 2 & Zeile 3 Spalte 3 & Zeile 3 Spalte 4\\
		Zeile 4 Spalte 1 &  Zeile 4 Spalte 2 & Zeile 4 Spalte 3 & Zeile 4 Spalte 4\\\bottomrule
	\end{tabular}
	\caption[Testtabelle]{\label{tab:tabelle1}Testtabelle}
\end{table}
\lipsum[1-2]

\begin{figure}
	\centering 
	\includegraphics{img/firmenlogo.jpg}
	\captionsetup{format=hang}
	\caption[Optionaler Kurztitel für das Abbildunggsverzeichnis]{\label{fig:test}Demo-Abbildung, um zu verdeutlichen, wie gleitende Objekte gesetzt werden und wie entsprechend die Quelle zitiert wird. \\Quelle: \cite[][S. 223]{TD15}}
\end{figure}
	
\subsection{Listings}	

Das Einbinden eines Listings mit der entsprechenden Umgebung ist auch kein Problem, wie man in Listing \vref{lst:helloworld} sehen kann. Schauen Sie sich hierzu das \texttt{listings}-Paket an! 
		
		\newpage
		
		
\lstset{language=Java}
\begin{lstlisting}[caption={Hello World!}, label={lst:helloworld}]
public static void main(String args[]) {
   System.out.println("Hello World!");
}
\end{lstlisting}


\subsection{Mathematische Formeln}
Auch mathematische Ausdrücke können mit \LaTeX~ sehr gut gesetzt werden, wie man anhand der Gleichungen \vref{eqn:e1} und \vref{eqn:e2} sehen kann -- konsultieren Sie hierzu bitte entsprechende Dokumentationen, die Online zur Verfügung stehen.
\begin{equation}
\left|{1\over N}\sum_{n=1}^N \gamma(u_n)-{1\over 2\pi}\int_0^{2\pi}\gamma(t){\rm d}t\right| \le {\varepsilon\over 3}.\\
\label{eqn:e1}
\end{equation}

\begin{equation}
f(x)=x^2
\label{eqn:e2}
\end{equation}


\subsection{Algorithmen}
Algorithmen können als Pseudocodes dargestellt und referenziert werden, wie z.\,B. in Algorithmus \vref{alg:euclid} -- sogar bis auf Zeilennummern
(siehe die \texttt{while}-Anweisung in Zeile \vref{alg:euclid:while}). Schauen Sie sich hierzu bitte das Paket \texttt{algorithmicx} an.



\begin{algorithm}
\begin{algorithmic}[1]
\Procedure{Euclid}{$a,b$}\Comment{The g.c.d. of a and b}
   \State $r\gets a\bmod b$
   \While{$r\not=0$}\Comment{We have the answer if r is 0} \label{alg:euclid:while}
      \State $a\gets b$
      \State $b\gets r$
      \State $r\gets a\bmod b$
   \EndWhile\label{euclidendwhile}
   \State \textbf{return} $b$\Comment{The gcd is b}
\EndProcedure
\end{algorithmic}
\caption{Euklid'scher Algorithmus}\label{alg:euclid}
\end{algorithm}


% !TEX root =  master.tex
\chapter{Einführungs}
\chapterauthor{Written by }


CovidNet beste \autocite{wang2020covid}.
\newpage
% !TEX root =  master.tex
\chapter{Weiterführung}



\newpage
% !TEX root =  master.tex
\chapter{Weiterführung}

\newpage
% !TEX root =  master.tex
\chapter{Conclusion}\label{chapter:conclusion}
\chapterauthor{Written by All}

In this project we wanted to create a \ac{CAD} like system that is capable of detecting and classifying COVID-19 infections based on \acl{CXR} images by solving the eponymous Kaeggle challenge. Our solution consists of a multi-stage approach including several deep neural networks to achieve the goal set by the challenge creators. We showed that a \ac{YOLO} and a Faster \ac{R-CNN} are suitable candidates for detecting and localizing COVID-19 regions of interests. Since the dataset required by the challenge is limited, we decided to include the \ac{RSNA} and \ac{NIH} data from previous challenges to pre-train our model weights.
We experienced that the Faster \ac{R-CNN} outperforms the \ac{YOLO} model on our test set, which is also consistent with other work in the field of object detection. However, for our final COVID-19 detection, we created an ensemble model that uses a weighted box fusion of both model outputs ultimately leading to the best results with a \ac{mAP}@$.50$ score of $0.451$. In general we are pleased with the final outcome of the COVID-19 detection performance and think that our work can be used as a solid basis for further research in this area. Due to the limited data we saw some overfitting effects with increased training time, which may be a starting point for further research. This may involve a repetition of our experiments with more data, an ablation study on how the pre-training of both models contribute to the performance or may be a root cause for overfitting effects, a hyper-parameter optimization process or even including more models with different architectures into the ensemble combination. Natural candidates would be for example a RetinaNet \autocite{lin_focal_2018} or a segmentation-based mask R-CNN \autocite{maskRcnn}.
Another interesting extension may include research about an end-to-end train process of the ensemble approach, meaning that there could be an improvement in performance if both single models would be trained together in a joint fashion. However, we only had limited compute resources and time budget which is why we leave this open for further research.

For the classification of COVID-19 infections we were not able to show that our \textit{ResNeXt} model could classify provided images into all the required classes with sufficient performance. However we were able to show that our model is able to decide between COVID-19 typical and negative cases to some extent. Although we were able to overcome the initial overfitting by training the classification layer only, the training did not produce stable and reliable results. The massive overfitting issues show that pre-training is an absolute requirement to solve the classification task of this challenge. To improve the performance either a more extensive pre-training on the NIH dataset could be done, although the pre-training on this dataset also showed already a slight overfitting, or additional pre-trainings on other datasets. There are additional large data sets of \ac{CXR} images publicly available, like the \textit{CheXpert} dataset by \citeauthor{irvin2019chexpert} \autocite{irvin2019chexpert}. On the other hand we could imagine that the limitation of the gradient updates to the classification layer somehow constrains or slows down the training of the model. Thus additional regularization methods should be tried like adding additional dropout layers to the hidden layers of the pre-trained models by redefining the forward path of the pre-trained network part.

In this report we have talked a lot about how different machine learning methods, and in this case computer vision, can help combat the pandemic. Honorably ML researchers and doctors all around the world scrambled to try and help at the start and during the pandemic by providing datasets, studies and models concerning all facets of what to do against the virus, for the infected or for society as a whole. Recently though, doubts have been cast on these efforts, as some studies suggest that machine learning models might not be very effective or in the worst case even dangerous. In a large scale meta study, published in the British Medical Journal, \citeauthor{wynants_prediction_2020} analyze 232 models proposed in 169 studies and find that not a single one is suitable for clinical use but that at least two might be promising with more work \autocite{wynants_prediction_2020}. Another study by \citeauthor{aix-covnet_common_2021} that closely examines 62 studies regarding the topic that we were also working on of identifying COVID-19 infections in patients, comes to the same conclusion that none of them are suitable for clinical use \autocite{aix-covnet_common_2021}. These concerns show that while there has been great effort and a lot of good will in the research community, the processes and studies are not yet as matured as maybe ML researchers and practitioners would like. While most times no harm is done in other areas of ML usage when an error is made, the stakes in the medical field are a lot higher and therefore a lot more care has to be put into every study and ML model that is proposed.

All in all, we liked working on this project. We think that with the global situation around COVID-19, ML may be a key-factor in solving different issues regarding medical image analysis while still keeping the aforementioned concerns in mind. Overall we are very pleased with our results with regards to the detection and localization of opacities. For the classification we expected better results. However we were able to provide solutions for all goals set by the underlying Kaggle Challenge. 

%\input{Kapitel/mitigation.tex}
%	Literaturverzeichnis
\newpage
\ihead{} % Neue Header-Definition
\printbibliography [title=References]


\cleardoublepage

% Der Anhang beginnt hier - jedes Kapitel wird alphabetisch aufgezählt. (Anhang A, B usw.)

\appendix
\ihead{\appendixname~\thechapter} % Neue Header-Definition
% appendix.tex einziehen
%\input{Kapitel/anhangtitle.tex}
%\chapter{Anhang}




% Ehrenwörtliche Erklärung ewerkl.tex einziehen
%% !TEX root =  master.tex

\clearpage
\chapter*{Ehrenwörtliche Erklärung}

% Wird die folgende Zeile auskommentiert, erscheint die ehrenwörtliche
% Erklärung im Inhaltsverzeichnis.

% \addcontentsline{toc}{chapter}{Ehrenwörtliche Erklärung}
Ich versichere hiermit, dass ich die vorliegende Arbeit
 mit dem Thema: \textit{\DerTitelDerArbeit} selbstständig verfasst und keine anderen als die angegebenen Quellen und
Hilfsmittel benutzt habe. Ich versichere zudem,
dass die eingereichte elektronische Fassung mit der gedruckten Fassung übereinstimmt.

\vspace{3cm}
Heidelberg, \today \hfill \DerAutorDerArbeit


\end{document}
